% !TEX program = xelatex
\documentclass[11pt,a4paper]{ctexart}

\usepackage{amsmath,amssymb,amsfonts}
\usepackage{booktabs}
\usepackage{geometry}
\usepackage{hyperref}
\usepackage{enumitem}
\usepackage{xcolor}
\usepackage{longtable}
\usepackage{array}

\geometry{left=2.5cm,right=2.5cm,top=2.5cm,bottom=2.5cm}
\hypersetup{colorlinks=true,linkcolor=blue,urlcolor=blue}

\newcolumntype{L}[1]{>{\raggedright\arraybackslash}p{#1}}

\title{%
  \textbf{星闪 SLB 物理层}\\[0.4em]
  \Large 6.2.1 一般描述技术文档\\[0.3em]
  \large 120\,kHz 内容分类 \textbullet\ 物理层信号上下文 \textbullet\ PIB 与两域结构
}
\author{依据 T/XS 10001-2025《星闪无线通信系统 接入层\\
同步低时延宽带空口 SLB 技术要求和测试方法》整理\\
（团体标准 20250326 版）}
\date{2026 年 7 月 6 日}

\begin{document}
\maketitle
\tableofcontents
\newpage

%======================================================================
\part{总览}
%======================================================================

\section{文档范围与模块划分}

本文档对应 SLB 120\,kHz 物理层协议第 6.2.1 节``一般描述''。该节在 6.1 物理层通用服务与 6.2.2 帧结构坐标系之上，\textbf{枚举 120\,kHz 子载波间隔系统空口传输的全部物理层内容类别}，并规定通信域与直通链路下\textbf{物理层信息块（PIB）}的组成边界。本文档为后续 6.2.2\textasciitilde 6.2.10 各子节的工程实现与联调提供统一索引。

6.2.1 本身不规定波形参数、帧结构或编码算法，而是定义三类内容集合（信号、控制信息、数据信息）及其在 PIB 中的打包方式。

\begin{table}[h]
\centering
\small
\begin{tabular}{L{3.2cm}L{4.8cm}L{5.5cm}}
\toprule
\textbf{逻辑子模块} & \textbf{子模块主题} & \textbf{核心输出/行为} \\
\midrule
6.2.1（信号） & 8 类物理层信号枚举 & 不承载高层比特的辅助波形分类；指向 6.2.3 波形实现 \\
6.2.1（控制） & 17 类物理层控制信息枚举 & 承载高层控制比特的分类；指向 6.2.4 字段定义 \\
6.2.1（数据） & 2 类物理层数据信息枚举 & 承载用户/信令载荷的分类；指向 6.2.5 处理流水线 \\
6.2.1（PIB） & 通信域 4 块 / 直通 2 块 & 信号+控制+数据的时频打包规则；定义散布资源与空闲 RE \\
\bottomrule
\end{tabular}
\end{table}

\section{与相关章节的关系}

\begin{table}[h]
\centering
\small
\begin{tabular}{L{2.2cm}L{2.2cm}L{9.1cm}}
\toprule
\textbf{相关节号} & \textbf{关系} & \textbf{说明} \\
\midrule
6.1 & 上游 & 物理层通用服务、G/T/附/直通链路命名、$T_s$/$f_s$ 与比特序约定 \\
6.2.2 & 上游 & 超帧/无线帧、RE 网格 $(k,l)$、五种 CP-OFDM 符号类型；提供``所在位置''坐标系 \\
6.2.3 & 下游 & 8 类信号的序列生成、RE 映射与功率（6.2.3.1\textasciitilde 10） \\
6.2.4 & 下游 & 17 类控制的比特字段、PIB 内布局与盲检规则 \\
6.2.5 & 下游 & 第一/二类数据的 TB 流水线、MCS 与 RE 映射 \\
6.2.6/6.5 & 下游 & 极化码/CRC/调制等跨模块基础能力 \\
6.2.7/6.2.8 & 下游 & SAB 盲检、FTS 捕获等系统流程，依赖本节内容分类 \\
6.3 & 平行 & 灵活带宽下相同内容分类的 $\alpha$ 参数缩放 \\
\bottomrule
\end{tabular}
\end{table}

\section{120\,kHz 系统定位与基本参数量}

\begin{table}[h]
\centering
\small
\begin{tabular}{lll}
\toprule
\textbf{概念} & \textbf{说明} & \textbf{标准条款} \\
\midrule
子载波间隔 & $\Delta f = f_s/256 = 120$\,kHz，$f_s = 30.72$\,MHz & 6.2.2.1 \\
资源元素 RE & 时频网格最小映射单位 $(k,l)$，$a_{80,l}\equiv 0$ & 6.2.2.7 \\
三类内容 & 信号（不承载高层比特）、控制（控制面）、数据（用户面/信令） & 6.1、6.2.1 \\
PIB 定义 & 为实现特定功能，映射多种信号/控制/数据组合的 RE 集合 & 6.2.1 \\
\bottomrule
\end{tabular}
\end{table}

\newpage
%======================================================================
\part{6.2.1 一般描述（工程解读）}
%======================================================================

\section{协议章节对应}

本节对应 T/XS 10001-2025 第 6.2 章第 6.2.1 节``一般描述''，涵盖：
\begin{itemize}[nosep]
  \item 8 类物理层信号；
  \item 17 类物理层控制信息；
  \item 2 类物理层数据信息；
  \item 通信域与直通链路 PIB 定义及散布资源枚举。
\end{itemize}
波形级实现分别见 6.2.3\textasciitilde 6.2.5；6.2.1 无独立编号子节。

\section{功能定义}

6.2.1 建立物理层\textbf{分类体系}与\textbf{信息块打包规则}，使接收端能够：
\begin{enumerate}[nosep]
  \item \textbf{分类识别}：区分 RE 上承载的是信号、控制还是数据；
  \item \textbf{块边界判定}：按 PIB 类型（SAB/RAB/G-PCIB/P-PCIB/A-PDIB）确定块边界与解码入口；
  \item \textbf{域类型区分}：识别通信域（经 G 节点，含 G/T/附链路）与直通链路（设备直连）的资源组织差异；
  \item \textbf{链路前缀解析}：按 G/T/附/直通链路前缀理解同名信号/控制在不同链路上的发送方向；
  \item \textbf{索引下游条款}：将每类内容映射至 6.2.3/6.2.4/6.2.5 的具体实现子节。
\end{enumerate}

\section{模块边界（上下游及职责划分）}

\subsection{上游模块}

\begin{table}[h]
\centering
\small
\begin{tabular}{L{4cm}L{4.5cm}L{5.5cm}}
\toprule
\textbf{上游} & \textbf{输入} & \textbf{说明} \\
\midrule
6.1 物理层一般描述 & 链路命名、$T_s$/$f_s$、比特序 & 全局 PHY 基准与三类内容定义 \\
MAC/RLC & 待发送控制/数据比特 & 由 6.2.4/6.2.5 组装后映射 \\
帧配置/调度 & 域类型、链路类型、符号类型 & 决定 PIB 选择与 RE 映射上下文 \\
6.2.2 帧结构 & 时频网格 $(k,l)_p$ & RE 坐标与符号编号 \\
\bottomrule
\end{tabular}
\end{table}

\subsection{下游模块}

\begin{table}[h]
\centering
\small
\begin{tabular}{L{4cm}L{4.5cm}L{5.5cm}}
\toprule
\textbf{下游} & \textbf{输出} & \textbf{说明} \\
\midrule
6.2.3 & 信号类型 + RE 需求 & FTS/STS、DMRS、SRS、CSI-RS、PMS、PAS \\
6.2.4 & 控制类型 + PIB 归属 & SAB/RAB/G-PCIB/P-PCIB 内 GCI/TCI \\
6.2.5 & 数据类型 + RE 映射 & 第一/二类 TB 与 DMRS-D 绑定 \\
6.2.6/6.5 & 待编码比特流 & CRC、极化码、调制、加扰 \\
6.2.7/6.2.8 & 搜网/同步流程入口 & 以 SAB/FTS/STS 为第一步 \\
\bottomrule
\end{tabular}
\end{table}

\subsection{本模块不负责的内容}

\begin{itemize}[nosep]
  \item CP-OFDM 波形与帧结构（6.2.2）；
  \item 各信号/控制/数据的比特字段与编码映射（6.2.3\textasciitilde 6.2.5）；
  \item 信道编码算法（6.2.6）、OFDM 合成（6.2.9）、射频发射（6.2.10）；
  \item 高层 MAC/RLC 语义与调度策略（数据链路层）。
\end{itemize}

\section{在 SLB 通信流程中的角色}

6.2.1 处于 PHY \textbf{内容组织最顶层}：在``帧结构已定''（6.2.2）与``具体波形生成''（6.2.3\textasciitilde 6.2.5）之间，回答``空口上传什么、放在哪个 PIB/链路/流程阶段''。

\textbf{通信域接收故事线（与 8 类信号出现顺序对应）：}
\begin{enumerate}[nosep]
  \item \textbf{搜网}：盲检 FTS/STS（同步信号）$\rightarrow$ 读 SAB 同步信息；
  \item \textbf{获知配置}：读广播信息（控制，散布资源）；
  \item \textbf{接入}：RAB 上 STS/FTS + 接入信息；
  \item \textbf{跟调度}：G-PCIB + G-DMRS-C $\rightarrow$ 解析 DSCI/预配置等；
  \item \textbf{传数据}：调度 RE 上的 DMRS-D + 第一/二类数据 $\pm$ 相位调整信号；
  \item \textbf{闭环}：CSI-RS/SRS 测量 $\rightarrow$ CQI；T-DMRS-C 伴随 TCI 反馈；
  \item \textbf{定位}：按资源指示在指定符号发/收位置测量信号；
  \item \textbf{空闲 RE}：未调度则不解析（$a_{k,l}{=}0$）。
\end{enumerate}

\textbf{直通链路简化故事线：}SAB $\rightarrow$ P-PCIB（直通 DSCI）$\rightarrow$ P-DMRS-C + DMRS-D + 数据。

%----------------------------------------------------------------------
\section{关键参数与强制要求（提炼版）}
%----------------------------------------------------------------------

\subsection{三类内容分类边界}

\begin{table}[h]
\centering
\small
\begin{tabular}{llll}
\toprule
\textbf{类型} & \textbf{条目数} & \textbf{是否承载高层比特} & \textbf{映射形态} \\
\midrule
物理层信号 & 8 & 否 & RE 集合；部分嵌入 SAB/RAB \\
物理层控制信息 & 17 & 是（控制面） & PIB 内或散布 RE \\
物理层数据信息 & 2 & 是（用户面/信令） & 调度 RE / A-PDIB \\
\bottomrule
\end{tabular}
\caption{三类内容维度（见 6.2.1、6.1）}
\end{table}

\subsection{8 类物理层信号索引}

\begin{table}[h]
\centering
\small
\begin{tabular}{clll}
\toprule
\textbf{序号} & \textbf{6.2.1 名称} & \textbf{典型链路} & \textbf{6.2.3 实现条款} \\
\midrule
1 & 同步信号 & G/直通/T & 6.2.3.1 FTS/STS \\
2 & G 链路公共解调参考信号 & G & 6.2.3.4 G-DMRS-C \\
3 & T 链路控制信息解调参考信号 & T & 6.2.3.7 T-DMRS-C \\
4 & 数据信息解调参考信号 & G/T/D/附/直通 & 6.2.3.6 DMRS-D \\
5 & 数据信息相位调整信号 & 数据链路 & 6.2.3.8 PAS \\
6 & 信道探测信号 & T$\rightarrow$G & 6.2.3.2 SRS \\
7 & 信道状态信息参考信号 & G$\rightarrow$T & 6.2.3.3 CSI-RS \\
8 & 位置测量信号 & G/T & 6.2.3.9 PMS \\
\bottomrule
\end{tabular}
\caption{8 类物理层信号（见 6.2.1）}
\end{table}

\subsection{17 类物理层控制信息索引}

\begin{table}[h]
\centering
\small
\begin{tabular}{cll}
\toprule
\textbf{序号} & \textbf{6.2.1 名称} & \textbf{6.2.4 实现子节} \\
\midrule
1 & 同步信息 & 6.2.4.1 \\
2 & 广播信息 & 6.2.4.2 \\
3 & 接入信息 & 6.2.4.4 \\
4 & 动态调度数据控制信息 & 6.2.4.6.1 \\
5 & 预配置调度激活去激活信息 & 6.2.4.6.2 \\
6 & 休眠与唤醒指示信息 & 6.2.4.6.3 \\
7 & 快速载波切换指示信息 & 6.2.4.6.4 \\
8 & 初始接入载波指示信息 & 6.2.4.6.5 \\
9 & G 链路位置测量信号或感知测量信号资源指示信息 & 6.2.4.6.6 \\
10 & T 链路位置测量信号资源指示信息 & 6.2.4.6.7 \\
11 & 直通链路动态调度数据控制信息 & 6.2.4.8 \\
12 & 单级 HARQ 反馈信息 & 6.2.4.10.1 \\
13 & 两级 HARQ 反馈信息 & 6.2.4.10.2 \\
14 & CQI 反馈信息 & 6.2.4.10.3 \\
15 & 调度请求信息 & 6.2.4.10.4 \\
\bottomrule
\end{tabular}
\caption{物理层控制信息索引（6.2.1 条文列举 15 项；含通信域/直通同步区分及 DSCI 多格式时，工程实现按 6.2.4 扩展为 17 类）}
\end{table}
\small 注：SAB/RAB/G-PCIB/P-PCIB 为 PIB 容器（6.2.4.2\textasciitilde 5、6.2.4.7、6.2.4.9），不单独计入控制信息清单。

\subsection{2 类物理层数据信息与 PIB 结构}

\begin{table}[h]
\centering
\small
\begin{tabular}{llll}
\toprule
\textbf{数据类型} & \textbf{调度特征} & \textbf{典型链路} & \textbf{标准条款} \\
\midrule
第一类物理层数据信息 & 仅预配置 & G/T/附 & 6.2.5.2 \\
第二类物理层数据信息 & 预配置 + 动态 & G/T/D/附 & 6.2.5.3 \\
\bottomrule
\end{tabular}
\end{table}

\textbf{通信域 PIB}（4 块）：SAB、RAB、G-PCIB、A-PDIB。

\textbf{通信域散布资源}（不进上述四块整包）：广播信息、G-DMRS-C、T 链路控制信息（TCI）、DMRS-D、物理层数据信息、信道探测信号、CSI-RS、位置测量信号、PAS、空闲 RE。

\textbf{直通链路 PIB}（2 块）：SAB、P-PCIB。

\textbf{直通链路散布资源}：P-DMRS-C、DMRS-D、物理层数据信息、空闲 RE。

%----------------------------------------------------------------------
\section{本节核心详解}
%----------------------------------------------------------------------

\subsection{物理层信号：功能与空口上下文总览}

6.2.1 列举 8 类\textbf{逻辑信号}；6.2.3 进一步细化为 FTS/STS、各类 DMRS 等\textbf{具体波形}。下表给出每类信号的功能定位、典型所在位置与流程上下文。

\begin{longtable}{@{}>{\small}L{2.4cm}>{\small}L{2cm}>{\small}L{3.2cm}>{\small}L{3.8cm}>{\small}L{3.2cm}@{}}
\caption{8 类物理层信号功能与空口上下文总览} \\
\toprule
\textbf{6.2.1 信号} & \textbf{6.2.3 条款} & \textbf{核心功能} & \textbf{典型时频/PIB 位置} & \textbf{流程上下文} \\
\midrule
\endfirsthead
\multicolumn{5}{c}{\small（续表）} \\
\toprule
\textbf{6.2.1 信号} & \textbf{6.2.3 条款} & \textbf{核心功能} & \textbf{典型时频/PIB 位置} & \textbf{流程上下文} \\
\midrule
\endhead
\bottomrule
\endfoot
同步信号 &
6.2.3.1 FTS/STS &
定时/频偏同步；STS $u$ 携带节点标识低 4 位 &
SAB/RAB：\#0=STS，\#1\textasciitilde\#2=FTS &
\textbf{空口第一步}：T 盲检 SAB（6.2.7）；多域 FTS 捕获（6.2.8） \\
G 链路公共解调参考信号 &
6.2.3.4 G-DMRS-C &
G 链路控制（G-PCIB）信道估计 &
含 SAB 的 TTI：首帧 \#6；无 SAB：\#0 &
G-PCIB 盲检/解调前；与 G 链路控制 RE 同功率 \\
T 链路控制信息解调参考信号 &
6.2.3.7 T-DMRS-C &
TCI（HARQ/CQI/SR）上行信道估计 &
TCI 发送符号的\textbf{第一个符号} &
DSCI/预配置触发 TCI 后；T$\rightarrow$G 反馈链路 \\
数据信息解调参考信号 &
6.2.3.6 DMRS-D &
第一/二类数据解调、MIMO 信道估计 &
调度数据符号内连续 $N$ 个 DMRS 符号 &
动态/预配置调度数据时；G/T/D/附/直通均可用 \\
数据信息相位调整信号 &
6.2.3.8 PAS &
跨符号相位跟踪/补偿 &
数据帧\textbf{首符号}上配置子载波 $k_0,k_1$ &
多符号数据、预编码/频偏导致相位漂移时 \\
信道探测信号 &
6.2.3.2 SRS &
T$\rightarrow$G 上行信道探测 &
T 链路符号；\texttt{srs-Set-Conf} 配置 &
G 测量 T 上行质量；辅助 T 链路调度 \\
信道状态信息参考信号 &
6.2.3.3 CSI-RS &
G$\rightarrow$T 下行信道参考 &
G 链路；CSI-RS 在 G-PCIB 后或 STS 后起始 &
T 测量下行质量 $\rightarrow$ CQI 反馈（6.2.4.10.3） \\
位置测量信号 &
6.2.3.9 PMS &
测距/定位（RTT、AoA 等） &
资源指示所调度 TTI 内指定符号 &
6.2.4.6.6/6.7 资源指示触发；按需插入 \\
\end{longtable}

\textbf{补充说明：}
\begin{itemize}[nosep]
  \item 直通链路 PIB 描述中另有 \textbf{P-DMRS-C}（6.2.3.5），功能类比 G-DMRS-C，固定映射于含 SAB 的 TTI 首帧 \#6，服务于 P-PCIB 解调；
  \item 6.2.3.10 定义\textbf{感知测量}参考信号，由 6.2.4.6.6 感知资源指示调度，实现感知时与 PMS 并列阅读。
\end{itemize}

\subsection{物理层控制信息：功能与典型上下文}

\begin{table}[h]
\centering
\small
\caption{17 类控制信息：功能与典型所在 PIB/上下文}
\begin{tabular}{L{4.5cm}L{5.5cm}L{4.5cm}}
\toprule
\textbf{控制信息} & \textbf{功能} & \textbf{典型位置/上下文} \\
\midrule
\multicolumn{3}{l}{\textit{网络发现与接入}} \\
同步信息 & 帧/节点/SAB 周期/GCI 位图 & SAB \#3\textasciitilde\#4 \\
广播信息 & 符号类型、TTI 长度、$N_{\mathrm{Symb-G-PCIB}}$ 等 & \#5、\#7；8\,ms 周期 \\
接入信息 & 随机接入、SR、服务发现等 & RAB \#3\textasciitilde\#4 \\
初始接入载波指示 & 三种带宽接入载波信道号 & G-PCIB；接入周期内 \\
\multicolumn{3}{l}{\textit{调度与资源}} \\
动态调度数据控制信息 & 时频资源、MCS、HARQ、CBG & G-PCIB \\
预配置调度激活去激活 & 半静态 grant 开/关 & G-PCIB \\
直通动态调度 & 直通 TB 时频资源 & P-PCIB \\
G/T 位置（感知）资源指示 & 测量符号/端口/波束配置 & G-PCIB \\
\multicolumn{3}{l}{\textit{功耗与载波}} \\
休眠与唤醒指示 & 58 T 节点 DRX 位图 & G-PCIB \\
快速载波切换指示 & 切换前后信道号与时刻 & G-PCIB \\
\multicolumn{3}{l}{\textit{反馈与请求}} \\
单级/两级 HARQ 反馈 & ACK/NACK、MCS 偏移 & TCI（T 链路） \\
CQI 反馈 & 各载波/子载波组 MCS 建议 & TCI \\
调度请求 & T 节点上行资源请求 & TCI 或 RAB 接入目的=1 \\
\bottomrule
\end{tabular}
\end{table}

\subsection{物理层数据信息（2 类）}

\begin{table}[h]
\centering
\small
\begin{tabular}{llll}
\toprule
\textbf{类型} & \textbf{调度方式} & \textbf{典型链路} & \textbf{配套信号/控制} \\
\midrule
第一类 & 仅预配置 & G/T/附 & 固定 $K{\cdot}N$；预配置激活 \\
第二类 & 预配置 + 动态 & G/T/D/附 & DSCI 指示资源/MCS/HARQ；DMRS-D \\
\bottomrule
\end{tabular}
\end{table}

\subsection{物理层信息块（PIB）与两域对比}

\begin{table}[h]
\centering
\small
\begin{tabular}{lll}
\toprule
\textbf{PIB} & \textbf{主要信号 + 控制} & \textbf{数据} \\
\midrule
SAB & FTS、STS + 同步信息 & --- \\
RAB & FTS、STS + 接入信息 & 接入目的=4 时 \#5\textasciitilde\#6 高层数据 \\
G-PCIB & G-DMRS-C（散布 \#6/\#0）+ 各类 GCI & --- \\
A-PDIB & STS、FTS、周期 DMRS-D & 附链路第一/二类数据 \\
P-PCIB（直通） & P-DMRS-C + 直通 DSCI & --- \\
\bottomrule
\end{tabular}
\end{table}

%----------------------------------------------------------------------
\section{数据类型、长度与维度}
%----------------------------------------------------------------------

6.2.1 为索引节，不规定各字段位宽；下表汇总\textbf{分类维度}与\textbf{条目规模}，供实现模块划分参考。具体位宽见 6.2.4/6.2.5。

\begin{table}[h]
\centering
\small
\begin{tabular}{llll}
\toprule
\textbf{对象/分类} & \textbf{条目数或规模} & \textbf{映射单位} & \textbf{备注} \\
\midrule
物理层信号类型 & 8 类 & RE 集合 & 不承载高层比特 \\
物理层控制类型 & 17 类 & PIB 内或散布 RE & 承载控制面比特 \\
物理层数据类型 & 2 类 & 调度 RE / A-PDIB & 含高层信令 \\
通信域 PIB & 4 块 & 多块 RE 组合 & SAB/RAB/G-PCIB/A-PDIB \\
直通链路 PIB & 2 块 & 多块 RE 组合 & SAB/P-PCIB \\
散布资源种类 & 9 类（通信域） & 独立 RE 映射 & 不进四块 PIB 整包 \\
空闲 RE & 未分配/未调度 & $a_{k,l}{=}0$ & 6.2.2.7 \\
\bottomrule
\end{tabular}
\end{table}

%----------------------------------------------------------------------
\section{时序关系与触发条件}
%----------------------------------------------------------------------

6.2.1 为静态分类，不含状态机；各类内容在超帧/无线帧内的\textbf{出现顺序与周期}由 6.2.2 帧结构与 6.2.4 调度规则决定。典型顺序如下：

\begin{enumerate}[nosep]
  \item \textbf{SAB}（含 FTS/STS + 同步信息）：周期出现，终端搜网入口；
  \item \textbf{广播信息}：8\,ms 周期，\#5/\#7 符号；
  \item \textbf{RAB}：接入阶段或 COT 末帧，竞争接入；
  \item \textbf{G-PCIB / P-PCIB}：每 TTI 或预配置周期，承载 GCI/直通调度；
  \item \textbf{数据 + DMRS-D}：按 DSCI/预配置映射，事件/周期触发；
  \item \textbf{SRS / CSI-RS / PMS}：预配置或资源指示触发，按需插入。
\end{enumerate}

%----------------------------------------------------------------------
\section{链路—信号—PIB 对照表}
%----------------------------------------------------------------------

\begin{table}[h]
\centering
\footnotesize
\begin{tabular}{L{3.2cm}L{2cm}L{2.2cm}L{2.2cm}L{4.5cm}}
\toprule
\textbf{内容} & \textbf{通信域 PIB} & \textbf{直通 PIB} & \textbf{链路} & \textbf{6.2.3/6.2.4 条款} \\
\midrule
FTS/STS + 同步信息 & SAB & SAB & G/直通 & 6.2.3.1、6.2.4.1\textasciitilde 2 \\
接入信息 & RAB & --- & T$\rightarrow$G & 6.2.4.4\textasciitilde 5 \\
G-DMRS-C & 散布 & --- & G & 6.2.3.4 \\
P-DMRS-C & --- & 散布 & 直通 & 6.2.3.5 \\
T-DMRS-C & --- & --- & T & 6.2.3.7 \\
DMRS-D & 散布 & 散布 & G/T/D/附/直通 & 6.2.3.6 \\
PAS & 散布 & 散布 & 数据链路 & 6.2.3.8 \\
SRS & 散布 & --- & T$\rightarrow$G & 6.2.3.2 \\
CSI-RS & 散布 & --- & G$\rightarrow$T & 6.2.3.3 \\
PMS & 散布 & --- & G/T & 6.2.3.9 \\
GCI（DSCI 等） & G-PCIB & --- & G & 6.2.4.6\textasciitilde 7 \\
直通 DSCI & --- & P-PCIB & 直通 & 6.2.4.8\textasciitilde 9 \\
TCI & --- & --- & T$\rightarrow$G & 6.2.4.10 \\
第一/二类数据 & 散布 / A-PDIB & 散布 & 各链路 & 6.2.5 \\
\bottomrule
\end{tabular}
\end{table}

%----------------------------------------------------------------------
\section{约束与实现要点}
%----------------------------------------------------------------------

\begin{enumerate}[nosep]
  \item \textbf{先同步后解调}：任何控制/数据解析均依赖 FTS/STS 建立的定时与 SAB 中的系统参数；
  \item \textbf{信号不承载用户比特}：8 类信号均为接收机辅助波形；用户载荷仅在 6.2.5 两类数据中；
  \item \textbf{链路前缀语义}：同名信号/控制按 G/T/附/直通前缀区分发送方向，实现时不得混用链路类型；
  \item \textbf{G-DMRS-C 载波约束}：无 SAB 的基础载波不发送 G-DMRS-C，盲检 G-PCIB 时需结合 SAB 存在性；
  \item \textbf{DMRS 与数据绑定}：DMRS-D 资源随 DSCI/预配置变化，不得按固定帧格式硬编码；
  \item \textbf{域类型不可混用}：通信域 PIB（RAB/G-PCIB/A-PDIB）不适用于直通链路；
  \item \textbf{空闲 RE 处理}：未分配 RE 对应 $a_{k,l}{=}0$，接收端不得当作有效载荷；
  \item \textbf{中心子载波}：$a_{80,l}\equiv 0$，任何内容均不映射至 \#80 子载波；
  \item \textbf{6.2.1 与 6.2.3 映射}：6.2.1 为逻辑分类，6.2.3.1\textasciitilde 10 为波形级定义（含感知扩展）；
  \item \textbf{PIB 与散布资源}：广播、DMRS-D、SRS、CSI-RS、PAS、PMS 等不进四块 PIB 整包，须单独解析映射规则。
\end{enumerate}

%----------------------------------------------------------------------
\section{与标准其他章节的衔接}
%----------------------------------------------------------------------

\begin{itemize}[nosep]
  \item \textbf{6.1}：物理层通用服务、三类内容定义、链路命名与 $T_s$/$f_s$ 基准；
  \item \textbf{6.2.2}：RE 网格、符号类型，是本文``所在位置''的坐标系；
  \item \textbf{6.2.3}：8 类信号序列、映射与功率的\textbf{实现细则}；
  \item \textbf{6.2.4}：17 类控制如何触发 SRS/CSI-RS/PMS 及数据调度；
  \item \textbf{6.2.5}：数据与 DMRS-D 的比特级处理；
  \item \textbf{6.2.7/6.2.8}：SAB 盲检与 FTS 捕获的系统流程；
  \item \textbf{6.3}：灵活带宽下相同内容分类的 $\alpha$ 缩放。
\end{itemize}

\vfill
\begin{center}
\small\textcolor{gray}{字段级定义与完整参数以 T/XS 10001-2025 第 6 章及 6.2.3\textasciitilde 6.2.5 各节为准；本文档提供 6.2.1 分类索引与空口上下文，供 SLB 120\,kHz 物理层实现与联调参考。}
\end{center}

\end{document}
